% Generated draft: requires mathematical and source review.
\documentclass[11pt]{article}
\usepackage[margin=25mm]{geometry}
\usepackage{fontspec}
\setmainfont{Latin Modern Roman}
\usepackage{amsmath,amssymb,mathtools,unicode-math}
\setmathfont{Latin Modern Math}
\usepackage{enumitem,needspace,xurl,hyperref}
\setlength{\parindent}{0pt}
\setlength{\parskip}{5pt}
\setlength{\emergencystretch}{3em}
\newcommand{\N}{\mathbb{N}}
\newcommand{\Z}{\mathbb{Z}}
\newcommand{\Q}{\mathbb{Q}}
\newcommand{\R}{\mathbb{R}}
\newcommand{\C}{\mathbb{C}}
\newcommand{\F}{\mathbb{F}}
\newcommand{\A}{\mathbb{A}}
\newcommand{\PP}{\mathbb{P}}
\newcommand{\E}{\mathbb{E}}
\newcommand{\eps}{\varepsilon}
\newcommand{\dd}{\,\mathrm{d}}
\begin{document}
\section*{UnsolvedMath ID: 30004598}
\textbf{Source problem code:} catalogue entry
\subsection{Definitions and mathematical statement}
For a fixed ground field \(K\), under the usual finite, reduced interpretation, a projective line arrangement is a finite set \(\mathcal A=\{L_1,\ldots,L_n\}\) of distinct projective lines in \(\PP^2_K\). The supplied source does not specify \(K\); this definition records a conventional interpretation rather than resolving the intended domain. Define the intersection set and incidence multiplicity by
\[
S(\mathcal A)=\{L_i\cap L_j:1\leq i<j\leq n\},
\qquad
m_{\mathcal A}(P)=\#\{i:P\in L_i\}.
\]
For the reduced union of these lines, the singular points are precisely the points of \(S(\mathcal A)\). A triple point is a point \(P\) with \(m_{\mathcal A}(P)=3\), so exactly three distinct arrangement lines pass through it.

The requested classification is of all such arrangements satisfying
\[
\forall P\in S(\mathcal A),\qquad m_{\mathcal A}(P)=3.
\]
Thus every intersection point must lie on exactly three arrangement lines; intersection points of multiplicity two or greater than three are excluded.

The classification target remains qualified: the source specifies neither the ground field and its characteristic nor the equivalence relation for classification. In particular, it does not state whether arrangements are to be classified up to projective transformations, by incidence structure, or in another sense. Its reference to the complex case occurs in the supplied literature assessment and does not explicitly restrict the problem statement to \(K=\C\). Finiteness and distinctness are conventional assumptions in the definition above, and the source gives no minimum number of lines or explicit nontriviality condition. If arrangements with fewer than two lines are admitted, the displayed condition holds vacuously.
\subsection{Short English statement}
Classify projective line arrangements in which every intersection point belongs to exactly three distinct lines. The intended field, classification equivalence, and treatment of trivial arrangements require clarification. The supplied record labels the problem open; this draft does not independently verify that status.
\subsection{Sources}
\begin{itemize}
\item Local UnsolvedMath registry: supplied record with permanent ID \(30004598\), source problem code \texttt{catalogue entry}, in the Oberwolfach Reports Open Problems collection. The record gives proposed year \(2021\) and names no individual proposer.
\item Supplied source citation: \emph{Logarithmic Vector Fields and Freeness of Divisors and Arrangements: New perspectives and applications} (2022), Oberwolfach Reports, pp.\ 225--300, DOI \texttt{10.4171/owr/2021/5}. \url{https://doi.org/10.4171/owr/2021/5}. The collection attributes the problem to workshop open-problem discussions documented in Oberwolfach Reports.
\item Supplied literature reference: Roulleau, Comptes Rendus Math\'ematique (2025). \url{https://doi.org/10.5802/crmath.741}.
\item Supplied literature reference, labeled ``Recent line-arrangement classification results.'' \url{https://doi.org/10.1007/s12215-024-01047-x}.
\item Supplied literature reference: Oberwolfach Report 5/2021. \url{https://doi.org/10.4171/owr/2021/5}.
\item Literature assessment dated exactly \texttt{2026-08-22}: the supplied curation labels the problem open and states that recent work adds infinite families of arrangements with only triple points while the complete classification remains open, including over \(\C\). This is supplied dated curation derived from the OWR record, not an independently verified literature assessment; its time-sensitive status and cited claims require independent checking.
\end{itemize}
\end{document}
